\chapter*{Abstract}
\addcontentsline{toc}{chapter}{Abstract}

% English Abstract
\section*{English}

Grapevine diseases pose a significant threat to global viticulture, causing substantial economic losses and compromising crop quality. Early and accurate disease detection is essential for effective disease management, yet traditional manual inspection methods are labor-intensive, subjective, and impractical for large-scale operations. While centralized cloud-based approaches using deep learning have demonstrated promising results, they require continuous network connectivity and introduce latency, making them unsuitable for remote vineyard environments.

This thesis presents the design and implementation of a fully autonomous edge AI system for real-time grape leaf disease detection deployed on resource-constrained ESP32-S3 microcontrollers. The system employs a two-stage deep learning pipeline: ESPDet-Pico, a YOLO-based lightweight object detection model, localizes individual grape leaves within captured images, followed by a fine-tuned MobileNetV2 classifier that identifies disease symptoms across four classes (Black Rot, Esca, Healthy, and Leaf Blight). To address the severe memory and computational constraints of embedded systems, both models undergo INT8 post-training quantization using the ESP-PPQ framework, reducing the detection model to 478 KB and the classification model to 2.73 MB while maintaining high accuracy.

A key innovation of this work is the hybrid weighted multi-leaf aggregation algorithm, which combines detection confidence, prediction entropy, and spatial quality metrics to synthesize robust disease predictions from multiple leaf observations within a single frame. This aggregation strategy achieves 98.7\% classification accuracy, representing a 3.4\% improvement over single-leaf predictions and reducing false positive rates by a factor of four.

The complete inference pipeline executes in 6.3 seconds per frame, including 78 ms for image capture, 366 ms for leaf detection, and 5.8 seconds for classifying up to ten detected leaves. Careful memory management enables operation within the ESP32-S3's 8 MB PSRAM, allocating buffers for camera DMA, JPEG frame storage, decoded RGB images, model activations, and processing intermediates. The system incorporates deep sleep power management with 5-minute duty cycles, achieving 20-25 days of autonomous operation on a 3000 mAh battery.

This research demonstrates the feasibility of deploying sophisticated deep learning models on ultra-low-cost microcontrollers for precision agriculture applications, offering a practical, scalable, and energy-efficient alternative to cloud-dependent solutions for vineyard disease monitoring.

\vspace{1cm}

\noindent\textbf{Keywords:} Edge AI, ESP32-S3, Deep Learning, Plant Disease Detection, YOLO, MobileNetV2, Model Quantization, Multi-Leaf Aggregation, Low-Power Computing, Precision Agriculture

\vspace{2cm}

% Italian Abstract
\section*{Riassunto}

Le malattie della vite rappresentano una minaccia significativa per la viticoltura globale, causando perdite economiche sostanziali e compromettendo la qualità delle colture. Il rilevamento precoce e accurato delle malattie è essenziale per una gestione efficace, tuttavia i metodi tradizionali di ispezione manuale sono laboriosi, soggettivi e impraticabili per operazioni su larga scala. Sebbene gli approcci centralizzati basati su cloud che utilizzano il deep learning abbiano mostrato risultati promettenti, richiedono connettività di rete continua e introducono latenza, rendendoli inadatti per ambienti viticoli remoti.

Questa tesi presenta la progettazione e l'implementazione di un sistema edge AI completamente autonomo per il rilevamento in tempo reale delle malattie delle foglie di vite, implementato su microcontrollori ESP32-S3 con risorse limitate. Il sistema utilizza una pipeline di deep learning a due stadi: ESPDet-Pico, un modello leggero di object detection basato su YOLO, localizza le singole foglie di vite all'interno delle immagini acquisite, seguito da un classificatore MobileNetV2 ottimizzato che identifica i sintomi delle malattie in quattro classi (Black Rot, Esca, Sana e Leaf Blight). Per affrontare i severi vincoli di memoria e calcolo dei sistemi embedded, entrambi i modelli subiscono una quantizzazione INT8 post-training utilizzando il framework ESP-PPQ, riducendo il modello di detection a 478 KB e il modello di classificazione a 2,73 MB mantenendo un'elevata accuratezza.

Un'innovazione chiave di questo lavoro è l'algoritmo di aggregazione multi-foglia con pesi ibridi, che combina la confidenza di detection, l'entropia della predizione e le metriche di qualità spaziale per sintetizzare previsioni robuste dalle osservazioni di più foglie in un singolo frame. Questa strategia di aggregazione raggiunge un'accuratezza di classificazione del 98,7\%, rappresentando un miglioramento del 3,4\% rispetto alle predizioni su singola foglia e riducendo i tassi di falsi positivi di un fattore quattro.

La pipeline di inferenza completa viene eseguita in 6,3 secondi per frame, includendo 78 ms per l'acquisizione dell'immagine, 366 ms per il rilevamento delle foglie e 5,8 secondi per la classificazione di fino a dieci foglie rilevate. Una gestione accurata della memoria consente l'operazione all'interno degli 8 MB di PSRAM dell'ESP32-S3, allocando buffer per il DMA della camera, la memorizzazione dei frame JPEG, le immagini RGB decodificate, le attivazioni del modello e gli intermedi di elaborazione. Il sistema incorpora una gestione energetica con deep sleep e cicli di lavoro di 5 minuti, raggiungendo 20-25 giorni di funzionamento autonomo con una batteria da 3000 mAh.

Questa ricerca dimostra la fattibilità di implementare modelli sofisticati di deep learning su microcontrollori a bassissimo costo per applicazioni di agricoltura di precisione, offrendo un'alternativa pratica, scalabile ed efficiente dal punto di vista energetico alle soluzioni dipendenti dal cloud per il monitoraggio delle malattie nei vigneti.

\vspace{1cm}

\noindent\textbf{Parole chiave:} Edge AI, ESP32-S3, Deep Learning, Rilevamento Malattie delle Piante, YOLO, MobileNetV2, Quantizzazione del Modello, Aggregazione Multi-Foglia, Computazione a Basso Consumo, Agricoltura di Precisione
